\subsection{Jerarquía de modelos base}
\label{SS:DES-MODEL-STRUCTURE}

Todas las entidades del dominio heredan de una jerarquía de tres clases
abstractas:

\begin{description}
  \item[\texttt{UUIDModel}] Proporciona un identificador primario
  \acs{uuid} v4. Frente a enteros autoincrementales, los \acs{uuid}
  evitan filtraciones de información (volumen de datos, orden de
  creación) y son seguros frente a colisiones en escalado horizontal.

  \item[\texttt{TimestampedModel}] Añade \texttt{created\_at} y
  \texttt{updated\_at} automáticos, proporcionando una traza de auditoría
  implícita.

  \item[\texttt{OrganizationOwnedModel}] Añade la clave foránea
  \texttt{organization} y dos \textit{managers}: \texttt{objects}
  (filtrado automáticamente por organización) y \texttt{unfiltered}
  (sin filtro, para tareas del sistema).
\end{description}

El modelo de datos se organiza en once subdominios, uno por aplicación
Django, con responsabilidades acotadas y dependencias explícitas. Los
diagramas de clase de cada subdominio se generan automáticamente con
\texttt{graph\_models} de \texttt{django-extensions} y se recogen en el
\cref{AP:CLASS-DIAGRAMS}; la vista global, con todas las entidades y
relaciones, se incluye en la \cref{fig:class-global}.


\subsection{Aislamiento multi-tenant}
\label{SS:DES-MODEL-ISOLATION}

El aislamiento entre organizaciones se garantiza mediante el
\texttt{TenantManager}, que filtra automáticamente cada \textit{query}
por el \texttt{organization\_id} del usuario autenticado. El contexto
del \dfn{tenant} se almacena en una variable de hilo y se establece
durante la autenticación \acs{jwt}, antes de que la vista procese la
petición. El flujo completo —desde el \texttt{POST /auth/login/} hasta
la limpieza del contexto por el \textit{middleware}— se diagrama en la
\cref{fig:seq-auth} (\cref{AP:SEQUENCE}).

El uso de \texttt{TenantManager} como \textit{manager} por defecto
hace estructuralmente imposible el olvido de un filtro de organización
en una \textit{query} manual, eliminando una fuente común de fugas de
datos entre \dfnpl{tenant}. La implementación concreta de los
\textit{managers} y del contexto de hilo se describe en
\cref{SS:IMPL-MULTITENANT}.


\subsection{Modelo de permisos contextuales}
\label{SS:DES-PERMISSIONS}

El sistema de autorización opera en tres capas
(\cref{fig:class-subjects}):

\begin{enumerate}
  \item El \textit{flag} \texttt{is\_staff} otorga permisos plenos de
  \dfn{gestor} sobre todas las asignaturas de la organización.

  \item La membresía de asignatura (\texttt{SubjectMembership}) define un
  rol (\texttt{COORDINATOR}, \texttt{TEACHER} o \texttt{STUDENT}) con
  trece permisos contextuales por defecto, configurables por organización
  y anulables individualmente por el \dfn{coordinador} mediante el campo
  \texttt{permission\_overrides}.

  \item Las reglas de asignación (\texttt{AssignmentRule}) restringen el
  acceso a instancias o problemas concretos.
\end{enumerate}

La resolución de permisos efectivos combina los valores por defecto del
rol, la configuración de la organización y las anulaciones del miembro,
con la restricción de que un \dfn{coordinador} no puede delegar un
permiso que él mismo no posee.
